\documentclass[spanish,a4paper,oneside,12pt,openright]{book}
\usepackage{graphicx}
\usepackage{amsmath,amsfonts,amssymb}
\usepackage[utf8]{inputenc}
\usepackage[T1]{fontenc}
%TNRoman:
%\usepackage{newtxtext,newtxmath}
%Arial-Helvetica:
%\usepackage[T1]{fontenc}
%\usepackage[scaled=0.92]{helvet}
%\renewcommand{\familydefault}{\sfdefault}
\usepackage[spanish, es-tabla]{babel}
\usepackage[Algoritmo]{algorithm}
\usepackage{algpseudocode}
\usepackage{url}
\usepackage{fancyhdr}
\usepackage{epstopdf}
\usepackage{minted}
\usepackage{csquotes}
\usepackage{rotating}
\usepackage{etoolbox}
\usepackage{zref-xr}
\usepackage{amssymb}
\usepackage{hyperref}
\usepackage[authoryear,round,longnamesfirst]{natbib}
\usepackage{multicol}
\usepackage{spverbatim}
\usepackage[pdf]{graphviz}
\usepackage{dirtree}
\usepackage{tabularx}
    \newcolumntype{L}{>{\raggedright\arraybackslash}X}
\usepackage{emptypage}
\usepackage[acronym,shortcuts,xindy]{glossaries}
\usepackage{markdown}
\usepackage{booktabs}
\usepackage{xcolor}
\usepackage{changepage}
\usepackage{tcolorbox}
\usepackage{listings}
\usepackage{datetime}

\newcommand{\capitalmonthname}[1]{%
  \ifcase#1\or
    Enero\or Febrero\or Marzo\or Abril\or Mayo\or Junio\or
    Julio\or Agosto\or Septiembre\or Octubre\or Noviembre\or Diciembre
  \fi
}

\newdateformat{monthyeardate}{%
  \capitalmonthname{\THEMONTH}, \THEYEAR}

% =============================================================================
% DATOS DE LA TESIS
% =============================================================================
\newcommand{\tfcTitle}{Diagramas de Bloques de Confiabilidad Extendidos con Colores: Formalización y Análisis}
\newcommand{\tfcAuthor}{Andrés Valentino Ferrero}
\newcommand{\tfcAdvisor}{Dr. Pedro R. D'Argenio}
\newcommand{\tfcGrade}{Licenciado en Ciencias de la Computación}
\newcommand{\tfcFaculty}{FACULTAD DE MATEMÁTICA, ASTRONOMÍA, FÍSICA Y COMPUTACIÓN}
\newcommand{\tfcCollege}{UNIVERSIDAD NACIONAL DE CÓRDOBA}

% Parche de compatibilidad para la página de certificación en main.tex
\newcommand{\authorname}{\tfcAuthor}

% Comandos auxiliares de edición
\newcommand{\citeme}[1]{{\color{red}{¡¡CITEME: #1!!}}}
\newcommand{\completar}[1]{{\color{red!75}{Completar: #1}}}

% =============================================================================
% MACROS MATEMÁTICAS Y DE NOTACIÓN (RBD / cRBD)
% =============================================================================
\newcommand{\pig}{\Pi_{G}}
\newcommand{\pigp}[1]{\Pi_{#1}}
\newcommand{\kutp}[1]{\mathcal{K}_{#1}}
\newcommand{\pigc}{{\Pi_{G}}_c}
\newcommand{\start}{\textsc{start}}
\newcommand{\nd}{\textsc{end}}
% Alias por si usás \snode y \tnode en las plantillas de capítulos
\newcommand{\rbd}{(G, \start, \nd)}
\newcommand{\rbdp}[1]{\ensuremath{(G_{#1}, \start_{#1}, \nd_{#1})}}
\newcommand{\graph}{G=(V,E)}
\newcommand{\graphp}[1]{\ensuremath{G_{#1}=(V_{#1},E_{#1})}}
\newcommand{\rgraph}{G=(V_f \uplus V_n,E)}
\newcommand{\walk}{v_1, ... , v_k}
\newcommand{\indgraph}[1]{G^{#1}}
\newcommand{\crbd}{(G, \start, \nd, \delta)}
\newcommand{\crbdp}[1]{\ensuremath{(G_{#1}, \start_{#1}, \nd_{#1}, \delta_{#1})}}

% Variables lógicas para la transformación de Tseitin
\newcommand{\reach}[2]{R_{#1, #2}} % Alcanzabilidad: \reach{n}{c} -> R_{n, c}
\newcommand{\fail}[1]{F_{#1}}      % Falla física:   \fail{n}     -> F_{n}

\setlength{\headheight}{15pt}

% wide page for side by side figures, tables, etc
\newlength{\offsetpage}
\setlength{\offsetpage}{1.0cm}
\newenvironment{widepage}{\begin{adjustwidth}{-\offsetpage}{-\offsetpage}%
    \addtolength{\textwidth}{2\offsetpage}}%
{\end{adjustwidth}}

\newtheorem{definition}{Definición}

\renewcommand{\listalgorithmname}{Índice de Algoritmos}

\addto\extrasspanish{
    \renewcommand{\sectionautorefname}{Sec.}
}

\makeglossaries

\DeclareMathOperator*{\argmax}{arg\,max}  % in your preamble
\DeclareMathOperator*{\argmin}{arg\,min}  % in your preamble

\DeclareMathOperator*{\minimize}{minimize}
\DeclareMathOperator*{\maximize}{maximize}


% =============================================================================
% DOCUMENTACIÓN DEL PAQUETE GLOSSARIES
% =============================================================================
% COMANDOS DISPONIBLES:
% \gls{RBD}      -> Primera vez muestra nombre completo + sigla; luego solo la sigla.
% \acrlong{RBD}  -> Siempre muestra la forma larga/expandida.
% \acrshort{RBD} -> Siempre muestra la forma corta/abreviada.
% \acrfull{RBD}  -> Siempre muestra la forma completa (larga + sigla).

% =============================================================================
% DEFINICIÓN DE ACRÓNIMOS DE LA TESIS
% =============================================================================
\newacronym{RBD}{RBD}{Diagrama de Bloques de Confiabilidad (\textit{Reliability Block Diagram})}
\newacronym{cRBD}{cRBD}{Diagrama de Bloques de Confiabilidad Coloreado (\textit{Colored Reliability Block Diagram})}
\newacronym{WMC}{WMC}{Conteo de Modelos Ponderado (\textit{Weighted Model Counting})}
\newacronym{CNF}{CNF}{Forma Normal Conjuntiva (\textit{Conjunctive Normal Form})}
\newacronym{WCNF}{WCNF}{Forma Normal Conjuntiva Ponderada (\textit{Weighted Conjunctive Normal Form})}
\newacronym{FT}{FT}{Árbol de Fallas (\textit{Fault Tree})}
\newacronym{FTs}{FTs}{Árboles de Falla (\textit{Fault Trees})}
\newacronym{DAG}{DAG}{Grafo Dirigido Acíclico (\textit{Directed Acyclic Graph})}

\lhead[]{\rightmark}
\chead[]{}
\rhead[\leftmark]{}
\renewcommand{\headrulewidth}{0.5pt}

\fancypagestyle{plain}{
\fancyhead[L]{}
\fancyhead[C]{}
\fancyhead[R]{}
\renewcommand{\headrulewidth}{0pt}
}

\pagestyle{fancy}

\widowpenalty 10000
\clubpenalty 10000

\renewcommand{\listingscaption}{Algoritmo}
\renewcommand{\listoflistingscaption}{Lista de algoritmos}

%\usemintedstyle{algol_nu}
%https://pygments.org/styles/

\setminted{
    breaklines=true
}

\renewcommand{\listingscaption}{Código}
\renewcommand{\thelisting}{\thesection.\arabic{listing}}
\preto\section{\setcounter{listing}{0}}

\begin{document}

% HOJA EN BLANCO DE GUARDA INICIAL
\pagenumbering{gobble}
\thispagestyle{empty}
\mbox{}
\newpage

% CARÁTULA
% Caratula según Anexo II - RHCD 572/2024 FAMAF
\begin{titlepage}

% LOGOS SUPERIORES
\noindent
\begin{minipage}[t]{0.5\linewidth}
    \flushleft
    \includegraphics[width=6cm]{00_tapa/figures/logo_unc.png}
\end{minipage}%
\begin{minipage}[t]{0.5\linewidth}
    \flushright
    \includegraphics[width=6cm]{00_tapa/figures/logo_famaf.png}
\end{minipage}

\vspace{2.5cm}

% TÍTULO Y AUTOR
\begin{center}
    {\LARGE \textbf{\tfcTitle}}\\[1.5cm]
    {\large por}\\[0.5cm]
    {\Large \textbf{\tfcAuthor}}
\end{center}

\vspace{1.5cm}

% TEXTO DESCRIPTIVO (Anexo II)
\noindent
Presentado ante la \textbf{\tfcFaculty} como parte de los requerimientos para la obtención del grado de \tfcGrade{} de la

\begin{center}
    \textbf{\tfcCollege}
\end{center}

\vspace{1cm}

\noindent
\begin{center}
    \monthyeardate\today
\end{center}

\vspace{0.5cm}

% DIRECTOR/A
\noindent
\begin{center}
Director: \tfcAdvisor
\end{center}

\vfill

% LICENCIA CREATIVE COMMONS (Texto oficial en castellano según Anexo II)
\begin{center}
    \includegraphics[width=3cm]{00_tapa/figures/ccby40.png}\\[0.3cm]
    Este trabajo se distribuye bajo una licencia
    \href{https://creativecommons.org/licenses/by/4.0/deed.es}{Creative Commons Atribución 4.0 Internacional}.
\end{center}

\end{titlepage} \cleardoublepage

\pagenumbering{gobble}
Certifico que el trabajo incluido en este documento es el resultado de tareas de investigación
originales y que no ha sido presentado para optar a un título en ninguna otra
Universidad o Institución.

\vspace{3cm}
\begin{flushright}
\tfcAuthor
\end{flushright}
\cleardoublepage

% RECONOCIMIENTOS + AGRADECIMIENTOS ============================================

%\chapter*{Reconocimientos}
%\pagenumbering{Roman} % para comenzar la numeracion de paginas en numeros romanos
%\addcontentsline{toc}{chapter}{Reconocimientos} % si queremos que aparezca en el índice
%\markboth{RECONOCIMIENTOS}{RECONOCIMIENTOS} % encabezado
%Esta Tesis fue desarrollada gracias a la ayuda económica de ...


%\phantom{p. 1}
\clearpage
\thispagestyle{empty}
\phantom{p. 2}
\clearpage

\chapter*{Agradecimientos} % si no queremos que añada la palabra "Capitulo"
\pagenumbering{Roman} % para comenzar la numeracion de paginas en numeros romanos
\addcontentsline{toc}{chapter}{Agradecimientos} % si queremos que aparezca en el índice
\markboth{AGRADECIMIENTOS}{AGRADECIMIENTOS} % encabezado

\completar{COMPLETAR}

% ABSTRACT ====================================================================

\phantom{p. 1}
\clearpage
\thispagestyle{empty}
\phantom{p. 2}
\clearpage

\chapter*{Resumen} % si no queremos que añada la palabra "Capitulo"
\addcontentsline{toc}{chapter}{Resumen} % si queremos que aparezca en el índice
\markboth{RESUMEN}{RESUMEN} % encabezado
\completar{COMPLETAR RESUMEN EN ESPAÑOL}

\phantom{p. 1}
\clearpage
\thispagestyle{empty}

\chapter*{Abstract} % si no queremos que añada la palabra "Capitulo"
\addcontentsline{toc}{chapter}{Abstract} % si queremos que aparezca en el índice
\markboth{ABSTRACT}{ABSTRACT} % encabezado
\completar{COMPLETING IN ENGLISH}


% \clearpage
% \phantom{p. 1}
% \thispagestyle{empty}

% GLOSARIO=====================================================================

% \renewcommand*\glspostdescription{\dotfill}
% \setlength{\glslistdottedwidth}{1\textwidth}
% \setlength{\glsdescwidth}{0.89\textwidth}
% \markboth{Glosario}{Glosario} % encabezado
\glsaddall
% \printglossary[type=\acronymtype,title=Glosario]
% \addcontentsline{toc}{chapter}{Glosario}


% CONTENIDOS ==================================================================
\phantom{p. 1}
\clearpage

\tableofcontents % indice de contenidos

\phantom{p. 1}
\clearpage
\thispagestyle{empty}
\phantom{p. 2}
\clearpage

% FIGURES =====================================================================

\cleardoublepage
\addcontentsline{toc}{chapter}{Índice de figuras} % para que aparezca en el indice de contenidos
\listoffigures % indice de figuras

\phantom{p. 1}
\clearpage
\thispagestyle{empty}
\phantom{p. 1}
\clearpage

% TABLES ======================================================================

% \cleardoublepage
% \addcontentsline{toc}{chapter}{Índice de tablas} % para que aparezca en el indice de contenidos
% \listoftables % indice de tablas

% pagina en blanco
\newpage

\cleardoublepage
\renewcommand\listoflistingscaption{Índice de códigos}
\listoflistings % Now typeset the list
\addcontentsline{toc}{chapter}{Índice de código} % para que aparezca en el indice de contenidos

% pagina en blanco
\newpage
\pagenumbering{arabic}
\phantom{p. 1}
\clearpage


% CONTENT =====================================================================

% Capítulo 1: Introducción
\chapter{Introducción}
\label{chapter:intro}

\section{Descripción general}
\label{sec:intro_descripcion}
\completar{Contexto general sobre confiabilidad de sistemas e importancia del modelado formal.}

\section{Problemática y motivación}
\label{sec:intro_problematica}
\completar{Limitaciones expresivas de los RBDs clásicos para sistemas multicapa/multiflujo y explosión combinatoria en la evaluación por caminos mínimos.}

\section{Contribución}
\label{sec:intro_contribucion}
\completar{Detalle de los aportes teóricos (formalización y codificación de Tseitin para cRBDs) y prácticos (herramienta en Python).}

\section{Organización del trabajo}
\label{sec:intro_overview}

El resto de este trabajo se estructura en dos partes principales:

En la \textbf{Parte~\ref{part:modelado}} se aborda el marco teórico y formal de los modelos gráficos de confiabilidad. En el Capítulo~\ref{chapter:rbd} se definen formalmente los \gls{RBD} como grafos dirigidos acíclicos, su semántica y sus operaciones de composición. Luego, en el Capítulo~\ref{chapter:crbd} se extiende este formalismo hacia los \gls{cRBD} sobre grafos dirigidos generales.

En la \textbf{Parte~\ref{part:codificacion_herramienta}} se presenta la resolución computacional y su implementación. El Capítulo~\ref{chapter:wmc} introduce el problema de \gls{WMC} y desarrolla tanto la codificación clásica por caminos como la traducción optimizada mediante la transformación de Tseitin. En el Capítulo~\ref{chapter:herramienta} se detalla la arquitectura, las funcionalidades y el modo de uso de la herramienta desarrollada. En el Capítulo~\ref{chapter:ejemplos} se evalúa la propuesta a través de distintos ejemplos y casos de estudio.

Finalmente, en el Capítulo~\ref{chapter:conclu} se exponen las conclusiones generales del trabajo y las líneas futuras de investigación. \cleardoublepage

% =============================================================================
\part{Extendiendo la expresividad de los Diagramas de Bloques de Confiabilidad}
\label{part:modelado}
% =============================================================================

% Capítulo 2: RBD
\chapter{Diagramas de Bloques de Confiabilidad (RBDs)}
\label{chapter:rbd}
En este capítulo se introduce el modelo matemático sobre el cual yace este trabajo. En \ref{sec:rbd_definicion} se establece la definición formal de \gls{RBD} (\textit{Reliability Block Diagram}). En \ref{sec:rbd_semantica} se describen definiciones de carácter semántico, como \textit{path} y \textit{cut set}. Luego, en \ref{sec:rbd_composiciones} se determinan los operadores del modelo: composición en serie y en paralelo.
Por último, se enuncia y demuestra una proposición que asegura que el modelo es composicional. Adicionalmente,
a lo largo del capítulo, se justifican algunas decisiones de diseño y restricciones importantes del modelo.
Posteriormente, en el siguiente capítulo, se definirá, en términos de RBDs, los RBD coloreados, que permiten extender la expresividad del modelo para que, además de composicional, sea cerrado bajo operadores y modular.

Previo al desarrollo de la definición formal de RBD, veamos los antecedentes teóricos necesarios.
\section{Preliminares}
\label{sec:rbd_prelim}




A \textit{directed graph} $G$ is a tuple $(V,E)$ where $V$ is a set of vertices and $E \subseteq V \times V$ is a set of ordered pairs. Given a graph $\graph$, a \textit{walk} is a sequence of vertices $\walk$ such that $\forall i: 1 \le i \le n-1 : (v_i, v_{i+1}) \in E$. If all vertices are distinct, a walk is called a \textit{path}. A \textit{cycle} is a walk where the only repeated vertices are $v_1=v_k$. A \textit{DAG (Directed Acyclic Graph)} is a directed graph that has no cycles.

Given a graph $\graph$ and a subset of nodes $X \subseteq V$, the \textit{induced graph} $\indgraph{X}$ is defined as: 
%
\[\indgraph{X} = (X, E \cap (X \times X)).\]



\section{RBD: Definición formal}
\label{sec:rbd_definicion}
\completar{Definición de RBD como grafo dirigido acíclico (DAG) con nodos terminales \start y \nd.}

\section{Semántica}
\label{sec:rbd_semantica}
\completar{Semántica de conectividad, estados de los componentes y función de estructura / probabilidad de supervivencia.}

\section{Composiciones}
\label{sec:rbd_composiciones}
\completar{Operaciones de composición (serie, paralelo, puente / grafos arbitrarios).} \cleardoublepage

% Capítulo 3: cRBD
\chapter{Diagramas de Bloques de Confiabilidad Coloreados (cRBD)}
\label{chapter:crbd}

\section{cRBD como grafo dirigido: Definición formal}
\label{chapter:crbd_definicion}
\completar{Definición de cRBD sobre grafos dirigidos (permitiendo ciclos/dependencias mutuas), alfabeto de colores y proyecciones.}

\section{Semántica}
\label{sec:crbd_semantica}
\completar{Semántica multicapa, compartición física de nodos entre distintos colores y condición de éxito del sistema.}

\section{Composiciones}
\label{sec:crbd_composiciones}
\completar{Operaciones de composición en cRBDs y acoplamiento entre canales.} \cleardoublepage

% =============================================================================
\part{Codificación Lógica, Implementación y Casos de Estudio}
\label{part:codificacion_herramienta}
% =============================================================================

% Capítulo 4: WMC y Codificación
\chapter{Codificación Lógica en Términos de WMC}
\label{chapter:wmc}

\section{Conteo de Modelos Ponderado (WMC)}
\label{sec:wmc_definicion}
\completar{Definición formal del problema SAT, CNF, WCNF y cálculo de probabilidades exactas mediante WMC.}

\section{Traducción basada en caminos}
\label{sec:wmc_caminos}
\completar{Codificación clásica mediante enumeración de Conjuntos Mínimos de Caminos (MPS) y el problema de la explosión combinatoria.}

\section{Traducción optimizada mediante transformación de Tseitin}
\label{sec:wmc_tseitin}
\completar{Codificación estructural directa empleando variables auxiliares de alcanzabilidad ($\reach{n}{c}$) y variables físicas de falla ($\fail{n}$).} \cleardoublepage

% Capítulo 5: La Herramienta
\chapter{Diseño e Implementación de la Herramienta}
\label{chapter:herramienta}

\section{Arquitectura general y flujo de procesamiento}
\label{sec:tool_arquitectura}
\completar{Visión general del pipeline en Python desde el archivo de entrada hasta el cálculo de confiabilidad.}

\section{Especificación y lectura de modelos}
\label{sec:tool_parser}
\completar{Formato de representación de la topología (nodos, aristas coloreadas, tasas de falla $\lambda$ y tiempo de misión $t$).}

\section{Generación de fórmulas lógicas y formato DIMACS WCNF}
\label{sec:tool_wcnf}
\completar{Implementación del método \texttt{to\_logic\_formula}, mapeo de identificadores de variables y escritura de cláusulas y pesos.}

\section{Integración con el solver GPMC y evaluación}
\label{sec:tool_gpmc}
\completar{Invocación del solver exacto GPMC, parseo de la salida y cómputo de la probabilidad final.}

\section{Visualización automática mediante Graphviz}
\label{sec:tool_graphviz}
\completar{Generación automática de diagramas topológicos en lenguaje DOT para RBDs y cRBDs.}

\section{Modo de uso}
\label{sec:tool_uso}
\completar{Requisitos del entorno, interfaz de línea de comandos (\texttt{python3 src/wmc\_crbd.py <filepath> <mission\_time>}) y formato de salida.} \cleardoublepage

% Capítulo 6: Ejemplos / Casos de Estudio
\chapter{Ejemplos y Casos de Estudio}
\label{chapter:ejemplos}

\completar{Introducción a los casos de prueba evaluados.}

\section{Sistemas con componentes compartidos y dependencias cruzadas}
\label{sec:ej_compartidos}
\completar{Ejemplos ilustrativos de cRBDs y comparación entre el enfoque por caminos y Tseitin.}

\section{Casos de estudio prácticos}
\label{sec:ej_casos_estudio}
\completar{Modelado de arquitecturas complejas (benchmarks FFORT, aviónica, sistemas de refrigeración nuclear, etc.).} \cleardoublepage

% =============================================================================
% Capítulo 7: Conclusiones
% =============================================================================
\chapter{Conclusiones y trabajo futuro}
\label{chapter:conclu}

\completar{CONCLUSIONES} \cleardoublepage


% APENDICES====================================================================


% \appendix



% BIBLIOGRAFIA ================================================================

\cleardoublepage
\addcontentsline{toc}{chapter}{Bibliografía}
\bibliographystyle{apalike}
\bibliography{tesis}

% HOJA FINAL DE AVAL DEL TRIBUNAL ======================
\cleardoublepage
\thispagestyle{empty}
\vspace*{3cm}
\noindent
Se deja constancia de que la presente versión final del Trabajo Especial titulado \textbf{``\tfcTitle''}, presentado por \textbf{\tfcAuthor}, se corresponde con la versión aprobada por el Tribunal Examinador.

\vspace{4cm}
\begin{center}
\begin{tabular}{ccc}
\rule{3.5cm}{0.4pt} & \rule{3.5cm}{0.4pt} & \rule{3.5cm}{0.4pt} \\[0.2cm]
Firma y Aclaración & Firma y Aclaración & Firma y Aclaración \\
Miembro del Tribunal & Miembro del Tribunal & Miembro del Tribunal
\end{tabular}
\end{center}

% HOJA EN BLANCO DE GUARDA FINAL
\clearpage
\thispagestyle{empty}
\mbox{}
\clearpage

\end{document}